% GENERATED FILE. Edit canonical lessons, then run npm run generate:books.
% canonical-source-hash: fnv1a64:b50c0f95666a9459
% canonical-lessons: FR-C498-escalade, FR-C498-grimper, FR-C498-gourde, FR-C498-lampe-de-poche, FR-C498-feu-de-camp

\chapter{En plein air}
\label{ch:fr-en-plein-air}
\hlchaptermodality{498}

\begin{chapteropening}
\textbf{By the end of this chapter:} \emph{I can talk about climbing and evenings outdoors.}

Complete the last lesson of chapter 498: I can talk about climbing and evenings outdoors.
\end{chapteropening}

\section[l'escalade]{l'escalade — climbing}

\label{lesson:FR-C498-escalade}



\begin{quote}
Before the new one: say the French for to fish, then the French for fisherman.
\end{quote}

\subsection*{You'll want to know: l'escalade}
\textbf{l'escalade} — \textquotedblleft{}climbing\textquotedblright{}.

\begin{quote}
Je fais de l'escalade en salle le mercredi.
\end{quote}

\begin{quote}
\emph{I go indoor climbing on Wednesdays.}
\end{quote}

Feminine. \emph{Faire de l'escalade} — to climb (as a sport).

\begin{cousinweb}
\emph{l'escalade} was borrowed from Italian \emph{scalata} (EtymDB), from \emph{scalare} \textquotedblleft{}to climb\textquotedblright{}, built on \emph{scala} \textquotedblleft{}ladder\textquotedblright{}, inherited from Latin \emph{scāla} (Lewis \& Short).
\end{cousinweb}

\begin{grammarlens}[title={What you've built}]
One more word for telling, explaining and describing at B1.
\end{grammarlens}

\subsection*{Your turn}
\begin{itemize}
\raggedright
  \item \emph{Say it:} \emph{l'escalade}
  \item \emph{Say it:} \emph{l'escalade}, once more
  \item \emph{Say it:} a sentence of your own with \emph{l'escalade}
  \item \emph{Recall:} say the French for to fish, then the French for fisherman, then say \emph{l'escalade} again
\end{itemize}

\begin{tcolorbox}[breakable,colback=teal!4,colframe=teal!35!black,title={Before you move on}]
What does l'escalade mean? (\textquotedblleft{}Climbing\textquotedblright{}.) Say it once more.
\end{tcolorbox}
\section[grimper]{grimper — to climb}

\label{lesson:FR-C498-grimper}



\begin{quote}
Before the new one: say the French for fisherman, then the French for climbing.
\end{quote}

\subsection*{You'll want to know: grimper}
\textbf{grimper} — \textquotedblleft{}to climb\textquotedblright{}.

\begin{quote}
Les enfants grimpent dans l'arbre.
\end{quote}

\begin{quote}
\emph{The children are climbing the tree.}
\end{quote}

\emph{Grimper à un arbre} or \emph{dans un arbre}. Also of prices: \emph{les prix grimpent}.

\begin{cousinweb}
\emph{grimper} is usually traced to \emph{gripper} \textquotedblleft{}to grab\textquotedblright{}, a verb of Germanic origin, with an added nasal sound; the origin is not settled.
\end{cousinweb}

\begin{grammarlens}[title={What you've built}]
One more word for telling, explaining and describing at B1.
\end{grammarlens}

\subsection*{Your turn}
\begin{itemize}
\raggedright
  \item \emph{Say it:} \emph{grimper}
  \item \emph{Say it:} \emph{grimper}, once more
  \item \emph{Say it:} a sentence of your own with \emph{grimper}
  \item \emph{Recall:} say the French for fisherman, then the French for climbing, then say \emph{grimper} again
\end{itemize}

\begin{tcolorbox}[breakable,colback=teal!4,colframe=teal!35!black,title={Before you move on}]
What does grimper mean? (\textquotedblleft{}To climb\textquotedblright{}.) Say it once more.
\end{tcolorbox}
\section[la gourde]{la gourde — water bottle, flask}

\label{lesson:FR-C498-gourde}



\begin{quote}
Before the new one: say the French for climbing, then the French for to climb.
\end{quote}

\subsection*{You'll want to know: la gourde}
\textbf{la gourde} — \textquotedblleft{}water bottle, flask\textquotedblright{}.

\begin{quote}
Avant de partir, il a rempli sa gourde.
\end{quote}

\begin{quote}
\emph{Before leaving, he filled his water bottle.}
\end{quote}

First a dried gourd. Informally, \emph{une gourde} is also a dimwit.

\begin{cousinweb}
\emph{la gourde} goes back to Old French \emph{coorde} (EtymDB), from Latin \emph{cucurbita} \textquotedblleft{}a gourd\textquotedblright{} (Lewis \& Short): dried gourds once served as flasks.
\end{cousinweb}

\begin{grammarlens}[title={What you've built}]
One more word for telling, explaining and describing at B1.
\end{grammarlens}

\subsection*{Your turn}
\begin{itemize}
\raggedright
  \item \emph{Say it:} \emph{la gourde}
  \item \emph{Say it:} \emph{la gourde}, once more
  \item \emph{Say it:} a sentence of your own with \emph{la gourde}
  \item \emph{Recall:} say the French for climbing, then the French for to climb, then say \emph{la gourde} again
\end{itemize}

\begin{tcolorbox}[breakable,colback=teal!4,colframe=teal!35!black,title={Before you move on}]
What does la gourde mean? (\textquotedblleft{}Water bottle, flask\textquotedblright{}.) Say it once more.
\end{tcolorbox}
\section[la lampe de poche]{la lampe de poche — torch, flashlight}

\label{lesson:FR-C498-lampe-de-poche}



\begin{quote}
Before the new one: say the French for to climb, then the French for water bottle.
\end{quote}

\subsection*{You'll want to know: la lampe de poche}
\textbf{la lampe de poche} — \textquotedblleft{}torch, flashlight\textquotedblright{}.

\begin{quote}
Il cherche ses clés avec une lampe de poche.
\end{quote}

\begin{quote}
\emph{He is looking for his keys with a torch.}
\end{quote}

Also \emph{une lampe torche}.

\begin{cousinweb}
\emph{la lampe de poche} is \textquotedblleft{}a pocket lamp\textquotedblright{}: \emph{lampe} comes from Latin \emph{lampas} \textquotedblleft{}a torch\textquotedblright{}, a borrowing from Greek, and \emph{poche} \textquotedblleft{}pocket\textquotedblright{} is of Frankish origin.
\end{cousinweb}

\begin{grammarlens}[title={What you've built}]
One more word for telling, explaining and describing at B1.
\end{grammarlens}

\subsection*{Your turn}
\begin{itemize}
\raggedright
  \item \emph{Say it:} \emph{la lampe de poche}
  \item \emph{Say it:} \emph{la lampe de poche}, once more
  \item \emph{Say it:} a sentence of your own with \emph{la lampe de poche}
  \item \emph{Recall:} say the French for to climb, then the French for water bottle, then say \emph{la lampe de poche} again
\end{itemize}

\begin{tcolorbox}[breakable,colback=teal!4,colframe=teal!35!black,title={Before you move on}]
What does la lampe de poche mean? (\textquotedblleft{}Torch, flashlight\textquotedblright{}.) Say it once more.
\end{tcolorbox}
\section[le feu de camp]{le feu de camp — campfire}

\label{lesson:FR-C498-feu-de-camp}



\begin{quote}
Before the new one: say the French for water bottle, then the French for torch.
\end{quote}

\subsection*{You'll want to know: le feu de camp}
\textbf{le feu de camp} — \textquotedblleft{}campfire\textquotedblright{}.

\begin{quote}
Le soir, nous chantions autour du feu de camp.
\end{quote}

\begin{quote}
\emph{In the evening we used to sing round the campfire.}
\end{quote}

\emph{Faire un feu de camp} — to make a campfire.

\begin{cousinweb}
\emph{le feu de camp} joins \emph{feu}, inherited from Latin \emph{focus} \textquotedblleft{}hearth\textquotedblright{}, and \emph{camp}, from Latin \emph{campus} \textquotedblleft{}a field\textquotedblright{} (Lewis \& Short).
\end{cousinweb}

\begin{grammarlens}[title={What you've built}]
One more word for telling, explaining and describing at B1.
\end{grammarlens}

\subsection*{Your turn}
\begin{itemize}
\raggedright
  \item \emph{Say it:} \emph{le feu de camp}
  \item \emph{Say it:} \emph{le feu de camp}, once more
  \item \emph{Say it:} a sentence of your own with \emph{le feu de camp}
  \item \emph{Recall:} say the French for water bottle, then the French for torch, then say \emph{le feu de camp} again
\end{itemize}

\begin{tcolorbox}[breakable,colback=teal!4,colframe=teal!35!black,title={Before you move on}]
What does le feu de camp mean? (\textquotedblleft{}Campfire\textquotedblright{}.) Say it once more.
\end{tcolorbox}
